Neural pre-decoders remove local errors from the syndrome of a quantum error correcting code before a global decoder
such as matching processes the remainder. We implement the two surface code pre-decoders that
NVIDIA released with pretrained weights, four and six layer three dimensional convolutional networks, as integer
networks on four AMD Alveo U55C FPGA cards, each card an independent node with two instances of one HLS kernel that
runs both networks. The kernel uses 8 bit products, 32 bit accumulation and a table based GELU, keeps all activations
on chip, and is bit exact with an integer model of the networks on every instance of every card. Measured with
NVIDIA's evaluation code, the six layer network lowers the logical error rate of matching by a factor of
\NumRatioAccFpThirteenHigh{} at physical error rate 0.006 and distance 13, as published, and the integer model by
\NumRatioAccIntThirteenHigh{}, while the four layer network gives little or no gain at
0.003. Four cards process \NumFourFast{} windows per second of the four layer network at distance 9 and
\NumFourAcc{} of the six layer network at distance 13, with linear scaling and within two percent of the kernel schedule,
but only \NumSpeedFast{} and \NumSpeedAcc{} times the same networks in PyTorch on
a 256 thread host CPU. A single window takes \NumLatFast{} and \NumLatAcc{} milliseconds, no faster than on the CPU. The arithmetic per window shows
that real time pre-decoding at microsecond syndrome cycles requires networks designed for hardware, while the cards
offload batch evaluation from the host.
